%% toc: P15 161 161.1,161.2,161.4,161.5
% Proofs for the comparison of almost-sure, probability, distribution and mean-power convergence.
\section{Almost-sure, probability, distribution and mean-power convergence: proofs}
\label{sec:annex-p15-w02b-p15-013-modes}

Use the definitions of Chapter~\ref{sec:p15-w02b-p15-013-modes}. All
random variables are finite and real, so differences \(X_n-X\) are
defined everywhere. Two facts are used repeatedly. For nonnegative
measurable \(V\), \(\mathbb EV=0\) exactly when \(V=0\) a.s.
(\cref{thm:p15-nonnegative-integral-laws}~(1)). For an interval
\(I\subseteq(0,1)\) and a constant \(c\geq0\), the variable
\(c\mathbf1_I\) on the uniform space is simple, and
\(\mathbb E(c\mathbf1_I)^r=c^r\,\mathbb U(I)\) for \(r>0\), where
\(\mathbb U(I)\) is the length of \(I\)
(\cref{def:p15-simple-integral,thm:p15-lebesgue-unit-interval}).

\begin{annexitem}{annex:p15-w02b-p15-013-power}{thm:p15-w02b-p15-013-power}{Concavity and convexity of powers, and the metrics}
\begin{proof}
(1) Let \(0<p\leq1\). For \(0<t\leq1\), \(t^{p-1}\geq1\) because
\(p-1\leq0\), so \(t^p=t\cdot t^{p-1}\geq t\); also \(0^p=0\). Hence
\(t^p\geq t\) on \([0,1]\). If \(a=b=0\) the first inequality is
\(0\leq0\). Otherwise put \(s=|a|+|b|>0\) and \(t=|a|/s\in[0,1]\), so
\(|b|/s=1-t\). Then
\[
 |a|^p+|b|^p=s^p\bigl(t^p+(1-t)^p\bigr)\geq s^p\bigl(t+(1-t)\bigr)
 =s^p\geq|a+b|^p,
\]
the last step because \(|a+b|\leq s\) and \(x\mapsto x^p\) is
increasing on \([0,\infty)\). For the second inequality, write
\(a=(a-b)+b\); the first gives \(|a|^p\leq|a-b|^p+|b|^p\). Exchanging
\(a\) and \(b\) gives \(|b|^p\leq|a-b|^p+|a|^p\). Together,
\(\bigl||a|^p-|b|^p\bigr|\leq|a-b|^p\).

(2) Let \(p\geq1\) and \(\varphi(x)=x^p\) on \([0,\infty)\). We show
\(\varphi((x+y)/2)\leq(\varphi(x)+\varphi(y))/2\) for \(x,y\geq0\).
By symmetry let \(x\leq y\), and put \(m=(x+y)/2\), \(h=(y-x)/2\). If
\(h=0\) there is nothing to prove. If \(h>0\), then \(m>x\geq0\), and
\(\varphi\) is continuous on \([x,y]\) and differentiable on \((x,y)\)
with \(\varphi'(\xi)=p\xi^{p-1}\). The mean value theorem gives
\(\eta\in(x,m)\) and \(\xi\in(m,y)\) with
\[
 \varphi(m)-\varphi(x)=p\eta^{p-1}h,\qquad
 \varphi(y)-\varphi(m)=p\xi^{p-1}h.
\]
Since \(0<\eta<\xi\) and \(p-1\geq0\), \(\eta^{p-1}\leq\xi^{p-1}\).
Hence \(\varphi(m)-\varphi(x)\leq\varphi(y)-\varphi(m)\), which is the
claim. Taking \(x=|a|\), \(y=|b|\),
\[
 |a+b|^p\leq(|a|+|b|)^p=2^p\Bigl(\frac{|a|+|b|}2\Bigr)^p
 \leq2^{p-1}\bigl(|a|^p+|b|^p\bigr).
\]

(3) Parts (1) and (2) give \(|a+b|^p\leq c_p(|a|^p+|b|^p)\), since
\(c_p=1\) for \(p\leq1\) and \(c_p=2^{p-1}\) for \(p\geq1\) (both
values are one at \(p=1\)). If \(X,Y\in L^p\) and \(c\) is real, then
\(\mathbb E|X+Y|^p\leq c_p(\mathbb E|X|^p+\mathbb E|Y|^p)<\infty\) and
\(\mathbb E|cX|^p=|c|^p\mathbb E|X|^p<\infty\), by monotonicity and
linearity of the nonnegative integral
(\cref{thm:p15-nonnegative-integral-laws}). Variables equal a.s.\ have
equal \(p\)th moments, and sums and multiples of a.s.\ equal variables
are a.s.\ equal, so the operations pass to classes. The vector-space
axioms hold pointwise.

(4) Let \(0<p<1\). By (3), \(d_p(X,Y)\) is finite on \(L^p\). It is
symmetric, and it is unchanged when \(X\) or \(Y\) is replaced by an
a.s.\ equal variable, so it is defined on classes. By the first fact
above, \(d_p(X,Y)=0\) exactly when \(|X-Y|^p=0\) a.s., that is, when
\(X=Y\) a.s. Pointwise, part (1) with \(a=X-Y\), \(b=Y-W\) gives
\(|X-W|^p\leq|X-Y|^p+|Y-W|^p\); integrating gives the triangle
inequality. For \(1\leq p<\infty\), \(\|\cdot\|_p\) is a norm on
\(L^p(\mathbb P)\) by \cref{thm:p15-w02b-p15-001-lp-results}~(3)--(4)
(and \cref{thm:p15-conditional-lp}), so \(\|X-Y\|_p\) is a metric.
Finally \(\|X_n-X\|_p\to0\) exactly when
\(\mathbb E|X_n-X|^p=\|X_n-X\|_p^p\to0\), because \(t\mapsto t^p\)
and \(t\mapsto t^{1/p}\) are continuous at zero on \([0,\infty)\) and
vanish there.

(5) By (3), pointwise
\(|X-Y|^p\leq c_p(|X-X_n|^p+|X_n-Y|^p)\), so
\(\mathbb E|X-Y|^p\leq c_p(\mathbb E|X_n-X|^p+\mathbb E|X_n-Y|^p)\to0\).
The left side does not depend on \(n\), so it is zero, and \(X=Y\) a.s.
by the first fact above.
\end{proof}
\end{annexitem}

\begin{annexitem}{annex:p15-w02b-p15-013-monotone}{thm:p15-w02b-p15-013-monotone}{H\"older with exponent \(p/q\)}
\begin{proof}
Let \(0<q<p<\infty\) and \(s=p/q>1\), with conjugate exponent
\(s'=s/(s-1)\in(1,\infty)\). \cref{thm:p15-holder-integral}, which
includes the infinite-norm cases, applied to \(f=|X|^q\) and \(g=1\),
gives
\[
 \mathbb E|X|^q\leq\bigl(\mathbb E|X|^{qs}\bigr)^{1/s}
                  \bigl(\mathbb E1^{s'}\bigr)^{1/s'}
 =\bigl(\mathbb E|X|^p\bigr)^{q/p},
\]
since \(\mathbb P(\Omega)=1\). Raising both sides to the power
\(1/q\), an increasing map of \([0,\infty]\), gives
\(\|X\|_q\leq\|X\|_p\). If \(X\in L^p\), the right side is finite,
so \(X\in L^q\). Applying the inequality to \(X_n-X\) gives
\(\mathbb E|X_n-X|^q\leq(\mathbb E|X_n-X|^p)^{q/p}\to0\) when
\(X_n\to X\) in \(L^p\). For \(p\geq1\), the case \(q=1\) (or \(p=1\)
itself) gives \(L^1\) convergence. If also \(X\in L^1\), then
\(\mathbb E|X_n-X|<\infty\) for all large \(n\), so \(X_n=(X_n-X)+X\)
is integrable for those \(n\), and
\(|\mathbb EX_n-\mathbb EX|\leq\mathbb E|X_n-X|\to0\) by
\cref{thm:p15-signed-integral-laws}.
\end{proof}
\end{annexitem}

\begin{annexitem}{annex:p15-w02b-p15-013-as-probability}{thm:p15-w02b-p15-013-as-probability}{Tail events, subsequences and the topological lemma}
\begin{proof}
(1) Let \(C\) be the convergence event, so \(\mathbb P(C)=1\). Fix
\(\varepsilon>0\) and put \(A_n=\{|X_n-X|>\varepsilon\}\). A point of
\(C\) lies in only finitely many \(A_n\), so
\(\limsup_nA_n\subseteq\Omega\setminus C\) and
\(\mathbb P(\limsup_nA_n)=0\). By
\cref{thm:p15-event-limit-calculus},
\(\limsup_n\mathbb P(A_n)\leq\mathbb P(\limsup_nA_n)=0\). Hence
\(\mathbb P(A_n)\to0\).

(2) Fix \(\varepsilon>0\). If \(|X-Y|>2\varepsilon\), then
\(|X_n-X|>\varepsilon\) or \(|X_n-Y|>\varepsilon\), by the triangle
inequality. Hence
\(\mathbb P(|X-Y|>2\varepsilon)\leq\mathbb P(|X_n-X|>\varepsilon)
+\mathbb P(|X_n-Y|>\varepsilon)\to0\), and the left side, which does
not depend on \(n\), is zero. Since
\(\{X\neq Y\}=\bigcup_{r\geq1}\{|X-Y|>2/r\}\), countable subadditivity
gives \(\mathbb P(X\neq Y)=0\).

(5) Suppose \(y_n\not\to y\). Then some open \(G\ni y\) is missed by
\(y_n\) for infinitely many \(n\); these indices form a subsequence
\(y_{n(m)}\notin G\). Every further subsequence stays outside \(G\) and
so does not converge to \(y\), contradicting the hypothesis.

(3) Suppose \(X_n\to X\) in probability, and let \(X_{n(m)}\) be a
subsequence. It also converges to \(X\) in probability, so by
\cref{thm:p15-probability-l1-tools}~(2) it has a further subsequence
converging to \(X\) a.s. Conversely, suppose every subsequence has a
further subsequence converging to \(X\) a.s. Fix \(\varepsilon>0\) and
let \(y_n=\mathbb P(|X_n-X|>\varepsilon)\), a real sequence. Given a
subsequence \(y_{n(m)}\), choose \(X_{n(m_k)}\to X\) a.s.; by (1),
\(y_{n(m_k)}\to0\). By (5) in \(\mathbb R\), \(y_n\to0\).

(4) Let \(X_{n(m)}\) be a subsequence. By (3) there is a further
subsequence with \(X_{n(m_k)}\to X\) on an event of probability one.
On that event, continuity of \(g\) gives
\(g(X_{n(m_k)})\to g(X)\). The variables \(g(X_n)\), \(g(X)\) are
measurable compositions with a continuous, hence Borel, map. So every
subsequence of \(g(X_n)\) has a further subsequence converging to
\(g(X)\) a.s., and the converse half of (3) gives
\(g(X_n)\to g(X)\) in probability.

(6) Suppose such a topology \(\tau\) exists, on the variables or on
their classes. Let \(T_n\) be the typewriter sequence of
\cref{thm:p15-w02b-p15-013-countermodels}~(2). Its proof below,
which does not use the present item, shows that
\(T_n\to0\) in probability and that \(T_n\) does not converge to \(0\)
a.s. By (3), every subsequence of \(T_n\) has a further subsequence
converging to \(0\) a.s., hence in \(\tau\). By (5), \(T_n\to0\) in
\(\tau\), hence \(T_n\to0\) a.s., a contradiction. Passing to classes
changes nothing, since a.s.\ convergence is insensitive to changing
each variable on a null set.

(7) Let \(N=\{\omega:\mathbb P(\{\omega\})=0\}\). It is a countable
union of null singletons, so \(\mathbb P(N)=0\). Let
\(\omega\notin N\), so \(c=\mathbb P(\{\omega\})>0\). If
\(X_n(\omega)\not\to X(\omega)\), there is \(\varepsilon>0\) with
\(|X_n(\omega)-X(\omega)|>\varepsilon\) for infinitely many \(n\), and
for each such \(n\), \(\mathbb P(|X_n-X|>\varepsilon)\geq c\). This
contradicts convergence in probability. Hence \(X_n\to X\) at every
point outside \(N\), that is, almost surely.
\end{proof}
\end{annexitem}

\begin{annexitem}{annex:p15-w02b-p15-013-comparison}{thm:p15-w02b-p15-013-comparison}{The implications}
\begin{proof}
(1) is \cref{thm:p15-w02b-p15-013-as-probability}~(1).

(2) Let \(\varepsilon>0\) and \(Z_n=X_n-X\). Apply
\cref{thm:p15-expectation-bounds} with \(\varphi(x)=|x|^p\), which is
nonnegative and Borel, and \(A=\{y:|y|\geq\varepsilon\}\), on which
\(\varphi\geq\varepsilon^p\):
\(\varepsilon^p\,\mathbb P(|Z_n|\geq\varepsilon)\leq\mathbb E|Z_n|^p\).
Hence \(\mathbb P(|X_n-X|>\varepsilon)\leq
\mathbb P(|X_n-X|\geq\varepsilon)\leq\varepsilon^{-p}\mathbb E|X_n-X|^p\to0\).
This is \cref{thm:p15-direct-variance-weak-laws}~(2) applied to \(Z_n\).

(3) is \cref{thm:p15-w02b-p15-013-monotone}.

(4) is \cref{thm:p15-weak-metric-premises}. There, convergence in
probability gives \(\alpha(X_n,X)\to0\), then \(\rho(F_n,F)\leq\alpha(X_n,X)\to0\),
hence \(F_n\Rightarrow F\); and \(X_n\Rightarrow c\) gives
\(\mathbb P(|X_n-c|>\varepsilon)\to0\).

(5) is \cref{thm:p15-probability-l1-tools}~(2).

(6) is \cref{thm:p15-w02b-p15-013-vitali}, including its two
sufficient conditions.

Compositions: almost-sure convergence and \(L^q\) or \(L^p\)
convergence each imply convergence in distribution through (1) or (2),
followed by (4). Every implication between two distinct modes among
the five named is therefore one of the eight
\[
 \begin{gathered}
 \text{a.s.}\Rightarrow\text{prob},\quad \text{a.s.}\Rightarrow\text{dist},\quad
 \text{prob}\Rightarrow\text{dist},\quad L^p\Rightarrow L^q,\\
 L^p\Rightarrow\text{prob},\quad L^p\Rightarrow\text{dist},\quad
 L^q\Rightarrow\text{prob},\quad L^q\Rightarrow\text{dist},
 \end{gathered}
\]
all proved, or one of the remaining twelve ordered pairs. Each of the
twelve fails in a model of \cref{thm:p15-w02b-p15-013-countermodels},
in which the hypothesis mode holds toward a limit and the conclusion mode
fails toward that limit (and, where stated, toward every limit):
\begin{itemize}
\item a.s.\(\Rightarrow L^p\), prob\(\Rightarrow L^p\),
      dist\(\Rightarrow L^p\), \(L^q\Rightarrow L^p\): model~(1),
      \(Z^{(p)}_n\to0\) a.s., in probability, in distribution and in
      \(L^q\), and in \(L^p\) toward no limit;
\item a.s.\(\Rightarrow L^q\), prob\(\Rightarrow L^q\),
      dist\(\Rightarrow L^q\): model~(1) with exponent \(q\),
      \(Z^{(q)}_n\);
\item \(L^p\Rightarrow\)a.s., \(L^q\Rightarrow\)a.s.,
      prob\(\Rightarrow\)a.s., dist\(\Rightarrow\)a.s.: model~(2),
      \(T_n\to0\) in every \(L^r\), in probability and in distribution,
      and a.s.\ toward no limit; model~(3) with \(p_n=1/n\) gives the
      same separations for independent variables;
\item dist\(\Rightarrow\)prob: model~(4).
\end{itemize}
That is \(4+3+4+1=12\) pairs. The models are defined for every
\(0<q<p<\infty\), so the list holds in each of the three ranges and for
\(q\) in any range below \(p\).
\end{proof}
\end{annexitem}

\begin{annexitem}{annex:p15-w02b-p15-013-vitali}{thm:p15-w02b-p15-013-vitali}{Reduction to the first mean through \(|X_n|^p\) and \(|X_n-X|^p\)}
\begin{proof}
Put \(W_n=|X_n|^p\), \(W=|X|^p\) and \(V_n=|X_n-X|^p\). These are
nonnegative, finite and measurable, and \(W_n\in L^1\) by hypothesis.
The map \(x\mapsto|x|^p\) is continuous, so
\cref{thm:p15-w02b-p15-013-as-probability}~(4) gives \(W_n\to W\) in
probability. Also \(V_n\to0\) in probability, because for \(\eta>0\),
\(\mathbb P(V_n>\eta)=\mathbb P(|X_n-X|>\eta^{1/p})\to0\). Uniform
integrability is tested with \cref{thm:p15-ui-small-sets}: a family is
uniformly integrable exactly when its first absolute moments are
bounded and it has uniform small-event control; every finite family of
integrable variables is uniformly integrable.

(1)\(\Leftrightarrow\)(3). Apply \cref{thm:p15-vitali-convergence} to
\(W_n\to W\) in probability. Its equivalence of uniform integrability
with \(\mathbb E|W_n|\to\mathbb E|W|<\infty\) is the present
equivalence, because \(|W_n|=|X_n|^p\) and \(|W|=|X|^p\).

(1)\(\Rightarrow\)(2). By (1)\(\Rightarrow\)(3),
\(\mathbb E|X|^p<\infty\), so \(X\in L^p\) and \(V_n\in L^1\) by
\cref{thm:p15-w02b-p15-013-power}~(3). Pointwise
\(V_n\leq c_p(W_n+W)\). Let \(K=\sup_n\mathbb EW_n<\infty\). Then
\(\sup_n\mathbb EV_n\leq c_p(K+\mathbb EW)<\infty\). Given
\(\epsilon>0\), small-event control of \((W_n)\) and of the single
integrable \(W\) give \(\delta>0\) such that \(\mathbb P(A)\leq\delta\)
implies \(\mathbb E(W_n;A)\leq\epsilon/(2c_p)\) for every \(n\) and
\(\mathbb E(W;A)\leq\epsilon/(2c_p)\); for such \(A\),
\(\mathbb E(V_n;A)\leq\epsilon\) for every \(n\). So \((V_n)\) is
uniformly integrable, and \cref{thm:p15-vitali-convergence}, applied to
\(V_n\to0\) in probability, gives \(\mathbb E|V_n-0|\to0\), that is,
\(\mathbb E|X_n-X|^p\to0\).

(2)\(\Rightarrow\)(1). Now \(X\in L^p\), each \(V_n\in L^1\), and
\(\mathbb EV_n\to0\); a convergent real sequence is bounded, so
\(\sup_n\mathbb EV_n<\infty\). Pointwise \(W_n\leq c_p(V_n+W)\), so
\(\sup_n\mathbb EW_n\leq c_p(\sup_n\mathbb EV_n+\mathbb EW)<\infty\).
Given \(\epsilon>0\), choose \(N\) with \(c_p\mathbb EV_n\leq\epsilon/2\)
for \(n>N\). The finite family \(W_1,\ldots,W_N,W\) has small-event
control, so there is \(\delta>0\) such that \(\mathbb P(A)\leq\delta\)
implies \(\mathbb E(W_n;A)\leq\epsilon\) for \(n\leq N\) and
\(c_p\mathbb E(W;A)\leq\epsilon/2\). For \(n>N\) and such \(A\),
\(\mathbb E(W_n;A)\leq c_p\mathbb E(V_n;A)+c_p\mathbb E(W;A)\leq
c_p\mathbb EV_n+\epsilon/2\leq\epsilon\). So \((W_n)\) is uniformly
integrable.

Sufficient conditions. If \(|X_n|\leq Z\) a.s.\ with \(Z\in L^p\),
then \(W_n\leq Z^p\) a.s.\ with \(Z^p\) integrable, and a dominated
family is uniformly integrable (\cref{thm:p15-ui-small-sets}). If
\(\sup_n\mathbb E|X_n|^r<\infty\) for some \(r>p\), let
\(\Phi(t)=t^{r/p}\) on \([0,\infty)\). It is finite, continuous and
hence Borel, and \(\Phi(t)/t=t^{r/p-1}\to\infty\) because \(r/p>1\).
Since \(\Phi(W_n)=|X_n|^r\), \cref{thm:p15-ui-moment-criteria} gives
uniform integrability of \((W_n)\). In both cases (1) holds and (2)
follows.
\end{proof}
\end{annexitem}

\begin{annexitem}{annex:p15-w02b-p15-013-countermodels}{thm:p15-w02b-p15-013-countermodels}{The five separating models}
\begin{proof}
(1) Let \(u\in(0,1)\). For \(n\geq1/u\), \(u\geq1/n\), so
\(Z^{(p)}_n(u)=0\); hence \(Z^{(p)}_n\to0\) at every point. By
\cref{thm:p15-w02b-p15-013-as-probability}~(1) and
\cref{thm:p15-w02b-p15-013-comparison}~(4), \(Z^{(p)}_n\to0\) in
probability and in distribution. For \(r>0\),
\(\mathbb E|Z^{(p)}_n|^r=n^{r/p}\,\mathbb U((0,1/n))=n^{r/p-1}\).
With \(r=q<p\) the exponent \(q/p-1\) is negative, so the moment tends
to zero; with \(r=p\) it equals one. Suppose \(Z^{(p)}_n\to X\) in
\(L^p\) for some \(X\). By
\cref{thm:p15-w02b-p15-013-comparison}~(2), \(Z^{(p)}_n\to X\) in
probability; it also tends to \(0\) in probability, so \(X=0\) a.s.\
by \cref{thm:p15-w02b-p15-013-as-probability}~(2). Then
\(\mathbb E|Z^{(p)}_n|^p\to0\), contradicting
\(\mathbb E|Z^{(p)}_n|^p=1\). Finally
\(|Z^{(p)}_n|^p=n\mathbf1_{(0,1/n)}\) has mean one, and for every
\(M>0\) and every \(n>M\),
\(\mathbb E\bigl(|Z^{(p)}_n|^p;|Z^{(p)}_n|^p>M\bigr)=1\). Thus the
supremum in the definition of uniform integrability is at least one
for every \(M\), and the family is not uniformly integrable.

(2) Each \(n\geq1\) has a unique representation \(n=2^k+j\),
\(0\leq j<2^k\): \(k\) is the largest integer with \(2^k\leq n\). Since
\(n<2^{k+1}\), \(k\to\infty\) as \(n\to\infty\). The set
\(I_n=[j2^{-k},(j+1)2^{-k})\cap(0,1)\) has length \(2^{-k}\) (for
\(k=0\) it is \((0,1)\)), so
\(\mathbb E|T_n|^p=\mathbb U(I_n)=2^{-k}\to0\) for every \(p>0\).
Hence \(T_n\to0\) in every \(L^p\), and by
\cref{thm:p15-w02b-p15-013-comparison}~(2) and~(4) in probability and
in distribution. Fix \(u\in(0,1)\) and \(k\geq1\). The \(2^k\) intervals
\([j2^{-k},(j+1)2^{-k})\), \(0\leq j<2^k\), are disjoint with union
\([0,1)\), so exactly one of them, say with \(j=j_k\), contains \(u\).
Then \(T_{2^k+j_k}(u)=1\), and since \(2^k\geq2\) there is
\(j\neq j_k\) with \(T_{2^k+j}(u)=0\). The indices \(2^k+j\) with
\(0\leq j<2^k\) lie in \([2^k,2^{k+1})\), disjoint blocks for distinct
\(k\). So \(T_n(u)=1\) for infinitely many \(n\) and \(T_n(u)=0\) for
infinitely many \(n\), and \(T_n(u)\) does not converge. The event on
which \(T_n\) converges to any finite \(X\) is therefore empty, and no
a.s.\ limit exists.

(3) Such independent variables exist on some probability space by
\cref{thm:p15-independent-sequence-existence}, with
\(\eta_n=(1-p_n)\delta_0+p_n\delta_1\). (i) For
\(0<\varepsilon<1\), \(\mathbb P(|B_n|>\varepsilon)=\mathbb P(B_n=1)=p_n\);
for \(\varepsilon\geq1\) it is zero. So \(B_n\to0\) in probability
exactly when \(p_n\to0\). Since \(B_n\in\{0,1\}\),
\(|B_n|^p=B_n\) and \(\mathbb E|B_n|^p=p_n\), so convergence in
\(L^p\) to zero also holds exactly when \(p_n\to0\). (ii) If
\(\sum_np_n<\infty\), the first Borel--Cantelli lemma
(\cref{thm:p15-borel-cantelli-first}) gives
\(\mathbb P(B_n=1\text{ i.o.})=0\); off that event \(B_n=0\) for all
large \(n\), so \(B_n\to0\) a.s. If \(\sum_np_n=\infty\), the events
\(\{B_n=1\}=B_n^{-1}(\{1\})\) are mutually independent because the
variables are, and the second Borel--Cantelli lemma
(\cref{thm:p15-borel-cantelli-second}) gives
\(\mathbb P(B_n=1\text{ i.o.})=1\); on that event \(B_n\not\to0\), so
\(B_n\) does not converge to zero a.s. For \(p_n=1/n\): \(p_n\to0\);
and \(\sum_n1/n=\infty\), because for every \(m\geq0\) the block
\(2^m<k\leq2^{m+1}\) has \(2^m\) terms each at least \(2^{-(m+1)}\),
so \(\sum_{k=1}^{2^M}1/k\geq1+M/2\) for every \(M\). Hence \(B_n\to0\)
in every \(L^p\) and in probability, and \(B_n=1\) infinitely often
a.s. If \(B_n\to X\) a.s.\ for some \(X\), then \(B_n\to X\) in
probability by \cref{thm:p15-w02b-p15-013-as-probability}~(1), so
\(X=0\) a.s.\ by part~(2) of that theorem, and \(B_n\to0\) a.s., which
has just been excluded.

(4) \(\mathbb P(\zeta=1)=\mathbb U((0,1/2))=1/2\) and
\(\mathbb P(\zeta=-1)=\mathbb U([1/2,1))=1/2\). The variable
\(-\zeta\) also takes the values \(\pm1\) with probability
\(1/2\) each, so every \(X_n\) has the same CDF as \(\zeta\), and
\(X_n\Rightarrow\zeta\) holds trivially. At every point
\(|X_n-\zeta|=2|\zeta|=2\). Hence
\(\mathbb P(|X_n-\zeta|>1)=1\) and
\(\mathbb E|X_n-\zeta|^p=2^p\) for every \(n\), and
\(X_n(u)-\zeta(u)=-2\zeta(u)\) does not tend to zero at any
\(u\).

(5) The intervals \(J_k=[2^{-k},2^{1-k})\), \(k\geq1\), are disjoint,
have union \((0,1)\) and have lengths \(2^{-k}\). So
\(Y=\sum_{k\geq1}2^k\mathbf1_{J_k}\) is a finite, positive,
measurable variable. By the countable additivity of the nonnegative
integral (\cref{thm:p15-nonnegative-integral-laws}),
\[
 \mathbb EY^p=\sum_{k\geq1}2^{kp}2^{-k}=\sum_{k\geq1}\bigl(2^{p-1}\bigr)^k.
\]
For \(0<p<1\) the ratio \(\varrho=2^{p-1}\) lies in \((0,1)\), and the
geometric series sums to \(\varrho/(1-\varrho)=1/(\varrho^{-1}-1)=
1/(2^{1-p}-1)\). For \(p=1\) every term equals one and the sum is
infinite, so \(\mathbb EY=\infty\). Since \(Y\) is finite,
\(Y(u)/n\to0\) at every \(u\). For \(0<p<1\),
\(\mathbb E|Y/n|^p=n^{-p}\,\mathbb EY^p\to0\). And
\(\mathbb E|Y/n|=\mathbb EY/n=\infty\) for every \(n\).

The allocation of the failing implications to the models, stated after
the list, is the content of (1)--(5) just proved; the count of pairs is
checked in \cref{annex:p15-w02b-p15-013-comparison}.
\end{proof}
\end{annexitem}

\begin{annexitem}{annex:p15-w02b-p15-013-finite-measure}{thm:p15-w02b-p15-013-finite-measure}{Lebesgue measure on the whole line}
\begin{proof}
(1) For fixed real \(x\), \(f_n(x)=0\) once \(n>x\), so \(f_n\to0\)
everywhere. The set \(\{|f_n|>1/2\}=[n,\infty)\) contains
\([n,n+m]\), of measure \(m\), for every integer \(m\geq1\)
(\cref{thm:p15-lebesgue-unit-interval}), so its measure is infinite.

(2) The intervals \([k,k+1)\), \(k\geq1\), are disjoint with measure
one, and \(f^2=\sum_kk^{-2}\mathbf1_{[k,k+1)}\). Countable additivity
of the nonnegative integral gives \(\int f^2\,d\lambda=\sum_kk^{-2}\)
and \(\int f\,d\lambda=\sum_kk^{-1}\). For \(k\geq2\),
\(k^{-2}\leq1/(k(k-1))=1/(k-1)-1/k\), and these differences telescope,
so \(\sum_{k=2}^Kk^{-2}\leq1-1/K<1\) and \(\int f^2\,d\lambda\leq2\).
The harmonic series diverges by the block estimate in the proof of
\cref{thm:p15-w02b-p15-013-countermodels}~(3). Hence
\(\|f\|_2\leq\sqrt2<\infty=\|f\|_1\), and
\(f\in L^2(\lambda)\setminus L^1(\lambda)\).
\end{proof}
\end{annexitem}
